
%----------------------------------------------------------------------------------------
%	CHAPTER 1 Introducción
%----------------------------------------------------------------------------------------

\chapterimage{chapter_head_2.pdf} % Chapter heading image

\chapter{Introducción}

La era de la mec\'{a}nica cu\'{a}ntica comenz\'{o} cuando Planck postul\'{o} que los modos de radiaci\'{o}n de frecuencia angular $\omega $ pod\'{\i}an contarse como part\'{\i}culas con energ\'{\i}a $E=\hbar \omega $, y deriv\'{o} la ley de radiaci\'{o}n de cuerpo negro sin llegar a infinitos. Las implicaciones f\'{\i}sicas de este postulado fundamental se hicieron m\'{a}s claras cuando Einstein mostr\'{o} que la cuantificaci\'{o}n de la radiaci\'{o}n tambi\'{e}n explicaba la dependencia de la frecuencia del efecto fotoel\'{e}ctrico. Esto pronto condujo a una proliferaci\'{o}n de nuevas ideas gracias a Bohr, Dirac, Born, Schr\"{o}dinger, Heisenberg y muchos otros, quienes aplicaron la idea de cuantificaci\'{o}n a part\'{\i}culas y encontraron un \'{e}xito notable en la descripci\'{o}n de fen\'{o}menos subat\'{o}micos. Dirac incluso construy\'{o} una teor\'{\i}a relativista del electr\'{o}n, que estaba en excelente acuerdo con los experimentos subat\'{o}micos. A pesar de su \'{e}xito, hubo algunas deficiencias esenciales de la mec\'{a}nica cu\'{a}ntica de part\'{\i}culas.

\section{\textquestiondown Porqu\'{e} una teoría Cu\'{a}ntica de Campos?}

Una de estas deficiencias fue filos\'{o}fica. La mec\'{a}nica cu\'{a}ntica comenz\'{o} con la descripci\'{o}n de la luz en t\'{e}rminos de fotones, pero la descripci\'{o}n cl\'{a}sica de la luz fue en t\'{e}rminos de propagaci\'{o}n de campos electromagn\'{e}ticos. Entonces, una teor\'{\i}a de los fotones requer\'{\i}a una receta de c\'{o}mo cuantificar campos, un eslab\'{o}n perdido entre las dos descripciones. Como se ver\'{a} en los pr\'{o}ximos cap\'{\i}tulos, la estructura necesaria para cuantificar los campos electromagn\'{e}ticos se puede utilizar para describir todas las part\'{\i}culas elementales como campos cu\'{a}nticos. Incluso se encontr\'{o} que la ecuaci\'{o}n de onda relativista de Dirac ten\'{\i}a la descripci\'{o}n m\'{a}s simple en t\'{e}rminos de campos cu\'{a}nticos que representan electrones y positrones.

Otro problema de la mec\'{a}nica cu\'{a}ntica de part\'{\i}culas es que, por definici\'{o}n, es v\'{a}lida en el r\'{e}gimen no relativista. Esto se debe s\'{o}lo a que utiliza hamiltonianos no relativistas para resolver varios problemas. De hecho, todo el dise\~{n}o de la mec\'{a}nica cu\'{a}ntica no relativista desaf\'{\i}a la relatividad. Por ejemplo, utiliza el concepto de potenciales, que es insostenible en cualquier teoría relativista ya que supone la transferencia de informaci\'{o}n a una velocidad infinita. Además, el espacio y el tiempo se tratan de manera muy diferente en la mecánica cu\'{a}ntica no relativista. Las coordenadas espaciales son operadores, mientras que el tiempo es un par\'{a}metro, y normalmente estudiamos la evoluci\'{o}n de diferentes operadores, incluidas las coordenadas espaciales, en el tiempo. La ecuaci\'{o}n de Dirac, aunque covariante, trata el espacio y el tiempo en bases diferentes. En una teoría verdaderamente relativista, el espacio y el tiempo deberían fusionarse en un espacio-tiempo, y no se puede hacer una distinci\'{o}n tan fundamental entre la parte espacial y la parte temporal.

Tambi\'{e}n hay un problema pr\'{a}ctico. En la naturaleza existen procesos en los que se crean o aniquilan nuevas part\'{\i}culas. Por ejemplo el decaimiento $\beta $ se puede considerar como la desintegraci\'{o}n de un neutr\'{o}n en un prot\'{o}n, un electr\'{o}n y un antineutrino:

\begin{equation}
n\rightarrow p+e^{-}+\overline{\nu }_{e}  \label{1}
\end{equation}

El neutr\'{o}n se aniquila en este proceso, mientras que se crean el prot\'{o}n, el electr\'{o}n, etc. La mec\'{a}nica cu\'{a}ntica de part\'{\i}culas se ocupa de las part\'{\i}culas estables y sus movimientos en varios potenciales. Sin cuantificaci\'{o}n de campos, cualquier c\'{a}lculo que implique la creaci\'{o}n o aniquilaci\'{o}n de part\'{\i}culas es totalmente err\'{o}neo. La teor\'{\i}a cu\'{a}ntica de campos, al incorporar la creaci\'{o}n y aniquilaci\'{o}n de cuantos o part\'{\i}culas como su caracter\'{\i}stica esencial, permiti\'{o} c\'{a}lculos significativos de resultados verificables experimentalmente. Cuando uno pasa al r\'{e}gimen relativista, se encuentra con la posibilidad de crear nuevas part\'{\i}culas a partir de la energ\'{\i}a. Por ejemplo, si un electr\'{o}n y un positr\'{o}n chocan con suficiente energ\'{\i}a, el proceso de colisi\'{o}n puede crear pares extra de electrones y positrones:

\begin{equation}
e^{-}+e^{+}\rightarrow e^{-}+e^{+}+e^{-}+e^{+}  \label{2}
\end{equation}

Esto, por supuesto, puede suceder s\'{o}lo si la energ\'{\i}a cin\'{e}tica total inicial del par $e^{-}~e^{+}$ es mayor que la energ\'{\i}a de masa en reposo del par adicional que se va a crear. Para permitir tales posibilidades, se requiere alg\'{u}n mecanismo para describir la creaci\'{o}n de part\'{\i}culas. Incluso cuando los estados inicial y final de un problema involucran las mismas part\'{\i}culas, un proceso cu\'{a}ntico puede pasar por estados intermedios que involucran a otras part\'{\i}culas que se crean y destruyen en el proceso. Tales estados intermedios pueden contribuir a la amplitud del proceso y afectar el resultado. Es importante considerar estas posibilidades cuando se quiere hacer pruebas de precisi\'{o}n de la teor\'{\i}a cu\'{a}ntica. La creaci\'{o}n y aniquilaci\'{o}n de las part\'{\i}culas intermedias en tales casos s\'{o}lo puede manejarse pasando a la teor\'{\i}a cu\'{a}ntica de campos.


\section{Operadores Aniquilador y Creador}

\newcommand{\Oaa}{\hat{a}}

Ahora, recordando un problema de la mec\'{a}nica cu\'{a}ntica no relativista que puede resolverse mediante la introducci\'{o}n de operadores que crean y aniquilan los cuantos, es el problema del oscilador arm\'{o}nico simple. Aqu\'{\i} se revisa este tema, ya que ser\'{a} muy \'{u}til para desarrollar la teor\'{\i}a cu\'{a}ntica de campos en los cap\'{\i}tulos siguientes.

El hamiltoniano cuantizado del oscilador unidimensional viene dado por,

\begin{equation}
\hat{H}=\frac{1}{2m}\left( \hat{P}^{2}+m^{2}\omega ^{2}\hat{X}^{2}\right) =%
\frac{\hbar \omega }{2}\left( \frac{\hat{P}^{2}}{m\hbar \omega }+\frac{%
m\omega }{\hbar }\hat{X}^{2}\right)  \label{3}
\end{equation}

\begin{exercise}
Demostrar que el Hamiltoniano asociado al oscilador armónico cuántico se lo puede escribir en términos de los operadores aniquilador y creador de la siguiente manera
\begin{equation*}
\hat{H}=\frac{\hbar \omega }{2}\left[ \Oaa^{\dagger }\Oaa+\Oaa \Oaa^{\dagger }\right] 
\end{equation*}
\end{exercise} 

\newcommand{\Op}{\hat{p}}
\newcommand{\Ox}{\hat{x}}
\newcommand{\OP}{\hat{P}}
\newcommand{\OX}{\hat{X}}

\begin{solution}
Se define las siguientes combinaciones de operadores adimensionales \emph{(Recordar que} $\hbar \omega$ \emph{\ posee unidades de energ\'{\i}a)}:

\begin{equation*}
\hat{p}=\frac{1}{\sqrt{m\hbar \omega }}\OP \qquad ;\qquad  \hat{x}=\sqrt{\frac{m\omega }{\hbar}}\OX
\end{equation*}

Reemplazando en la ecuaci\'{o}n (\ref{3}) asociada a $\hat{H}$

\begin{equation}
\hat{H}=\frac{\hbar \omega }{2}\left( \Op^{2}+\Ox^{2}\right)  \label{1.3}
\end{equation}

Con el fin de redefinir el Hamiltoniano de la l\'{\i}nea (\ref{1.3}), definir los siguientes operadores

\begin{equation}
\Oaa \equiv \frac{1}{\sqrt{2}}\left( \Op-\ii \Ox \right) \qquad ;\qquad \Oaa^{\dagger } \equiv \frac{1}{\sqrt{2}}\left( \Op+\ii \Ox\right)  \label{1.3A}
\end{equation}

que escrito en t\'{e}rminos de los operadores $\OX$ y \thinspace $\OP$ se reducen a

\begin{eqnarray*}
\Oaa &=&\frac{1}{\sqrt{2}}\left( \Op-\ii \Ox \right) =\frac{1}{\sqrt{2}}\left( \frac{1}{\sqrt{m\hbar \omega }} \OP- \ii \sqrt{\frac{m\omega }{\hbar }}\OX \right) =\frac{1}{\sqrt{2m\hbar \omega }}\left( \OP-\ii \sqrt{\frac{m\omega}{\hbar}}\sqrt{m\hbar \omega} \OX \right) \\
&=&\frac{1}{\sqrt{2m\hbar \omega }}\left( \OP-\ii m\omega \OX \right),
\end{eqnarray*}

en tanto que

\begin{eqnarray*}
\Oaa^{\dagger } &=& \frac{1}{\sqrt{2}}\left( \Op+ \ii \Ox\right) =\frac{1}{\sqrt{2}}\left( \frac{1}{\sqrt{m\hbar \omega }} \OP + \ii \sqrt{\frac{m\omega }{\hbar}} \OX \right) =\frac{1}{\sqrt{2m\hbar \omega }}\left( \OP+\ii \sqrt{\frac{m\omega}{\hbar}}\sqrt{m\hbar \omega } \OX \right) \\
&=&\frac{1}{\sqrt{2m\hbar \omega }}\left( \OP+\ii m\omega \OX \right), 
\end{eqnarray*}

concluyendo
\begin{equation}
\boxed{
\Oaa \equiv \frac{1}{\sqrt{2m\hbar \omega }}\left( \OP-\ii m\omega \OX \right) \qquad ;\qquad \Oaa^{\dagger }=\frac{1}{\sqrt{2m\hbar \omega }}\left( \OP+\ii m\omega \OX\right).}  \label{1.4}
\end{equation}

\noindent
Por otro lado, es necesario calcular las siguientes relaciones a partir de (\ref{1.3A}):

\vspace{0.2cm}
\noindent
\begin{itemize}
\item[$i)$]
\begin{eqnarray}
\Oaa+\Oaa^{\dagger } &=& \frac{1}{\sqrt{2}}\left( \Op+\ii \Ox \right) +\frac{1}{\sqrt{2}}\left( \Op-\ii \Ox \right) =\frac{1}{\sqrt{2}}\left( \Op+\ii \Ox +\Op - \ii \Ox \right) =\frac{2p}{\sqrt{2}}  \notag \\
&\therefore &~\ \Op=\frac{\sqrt{2}}{2}\left( \Oaa+\Oaa^{\dagger }\right) \quad \text{ó} \quad  \OP =  \sqrt{\frac{m \hbar \omega}{2}} \left( \Oaa+\Oaa^{\dagger} \right).
\label{1.4a}
\end{eqnarray}

\item[$ii)$]
\begin{eqnarray}
\Oaa-\textbf{}a^{\dagger } &=&\frac{1}{\sqrt{2}}\left( \Op+\ii \Ox \right) -\frac{1}{\sqrt{2}}\left( \Op-\ii \Ox \right) =\frac{1}{\sqrt{2}}\left( \Op+\ii \Ox-\Op+\ii \Ox\right) =\frac{2ix}{\sqrt{2}}  \notag \\
&\therefore &~\ \Ox=\frac{\sqrt{2}}{2\ii}\left( \Oaa-\Oaa^{\dagger }\right) \quad \text{ó} \quad  \OX = -\ii \sqrt{\frac{\hbar}{2m \omega}} \left( \Oaa-\Oaa^{\dagger }\right). 
\label{1.4b}
\end{eqnarray}
\end{itemize}

Empleando los resultados (\ref{1.4a}) y (\ref{1.4b}) derivados anteriormente en la relaci\'{o}n (\ref{1.3}),

\begin{eqnarray*}
\hat{H} &=&\frac{\hbar \omega }{2}\left( \Op^{2}+\Ox^{2}\right) =\frac{\hbar \omega }{2}\left[ \left( \frac{\sqrt{2}}{2}\left( \Oaa+\Oaa^{\dagger }\right) \right)^{2}+\left( \frac{\sqrt{2}}{2\ii}\left( \Oaa-\Oaa^{\dagger }\right) \right) ^{2} \right] \\
%%
&=&\frac{\hbar \omega }{2}\left[ \left( \frac{\sqrt{2}}{2}\right) ^{2}\left(\Oaa+\Oaa^{\dagger }\right) ^{2}+\left( \frac{\sqrt{2}}{2 \ii}\right) ^{2}\left(\Oaa-\Oaa^{\dagger }\right) ^{2}\right] =\frac{\hbar \omega }{2}\left[ \frac{2}{4}\left( \Oaa+\Oaa^{\dagger }\right) ^{2}+\frac{2}{4(-1)}\left( \Oaa-\Oaa^{\dagger}\right) ^{2}\right] \\
%&=&\frac{\hbar \omega }{4}\left[ \left( a+a^{\dagger }\right) \left(a+a^{\dagger }\right) -\left( a-a^{\dagger }\right) \left( a-a^{\dagger}\right) \right] =\frac{\hbar \omega }{4}\left[ aa+aa^{\dagger }+a^{\dagger}a+a^{\dagger }a^{\dagger }-\left( aa-aa^{\dagger }-a^{\dagger }a+a^{\dagger}a^{\dagger }\right) \right] \\
%&=&\frac{\hbar \omega }{4}\left[ aa+aa^{\dagger }+a^{\dagger }a+a^{\dagger}a^{\dagger }-aa+aa^{\dagger }+a^{\dagger }a-a^{\dagger }a^{\dagger }\right]=\frac{\hbar \omega }{4}\left[ aa^{\dagger }+a^{\dagger }a+aa^{\dagger}+a^{\dagger }a\right] \\
&=&\frac{\hbar \omega }{4}\left[ 2\Oaa \Oaa^{\dagger }+2\Oaa^{\dagger }\Oaa\right] =\frac{\hbar \omega }{2}\left[ \Oaa \Oaa^{\dagger }+\Oaa^{\dagger } \Oaa\right],
\end{eqnarray*}

demostrando que,

\begin{equation}
\boxed{
\hat{H}=\frac{\hbar \omega}{2}\left[ \Oaa^{\dagger }\Oaa+\Oaa \Oaa^{\dagger }\right].
\label{1.5}
}
\end{equation}
 \end{solution}

\vspace{0.3cm}
\noindent
Observe que si se estuviera en el caso cl\'{a}sico asociado $p$ y $x$, se puede tener escrito $\hat{H}=\hbar \omega \Oaa^{\dagger } \Oaa=\hbar \omega \Oaa \Oaa^{\dagger }$, sin embargo en el caso cu\'{a}ntico esto no es posible dado que $x\rightarrow \OX$ y $p\rightarrow \OP$, representan los operadores de ``posición'' y ``momento'' que satisfacen relaci\'{o}n de conmutaci\'{o}n:

\begin{equation}
\left[ \OX_{i},\OP_{j}\right] =i\hbar \delta _{ij}~\ \ \ \implies ~\ \ \ \left[\OX,\OP\right] =i\hbar,  \label{1.6}
\end{equation}

donde
\begin{equation*}
\left[ \OX,\OP\right] =\OX \OP-\OP \OX.
\end{equation*}

\noindent
Las relaciones anteriores conllevan a que las cantidades $\Oaa$\ y $\Oaa^{\dagger }$, definidas en la ecuaci\'{o}n (\ref{1.4}), tambi\'{e}n sean operadores. Cabe resaltar que a diferencia de los operadores $\OX$ y $\OP$, $\Oaa$\ y $\Oaa^{\dagger }$ no son operadores hermiticos, además satisfacen las relaciones de conmutación:

\begin{eqnarray}
\left[ \Oaa,\Oaa^{\dagger }\right] &=&1, \label{1.7} \\
\left[ \Oaa^{\dagger },\Oaa^{\dagger }\right] &=& 0 , \label{1.8} \\
\comm\big{\Oaa}{\Oaa} &=& 0 . \label{1.8aa}
\end{eqnarray}

\begin{exercise}
 Demostrar las relaciones (\ref{1.7}) a (\ref{1.8aa})
\end{exercise}

\begin{solution}
\vspace{1mm}
Haciendo uso de las expresiones (\ref{1.4}) en (\ref{1.7}):


\begin{eqnarray*}
\left[ \Oaa, \Oaa^{\dagger }\right] &=&\left[ \frac{1}{\sqrt{2m\hbar \omega }} \left( \OP - \ii m \omega \OX \right) ,\frac{1}{\sqrt{2m\hbar \omega }}\left( \OP+\ii m\omega \OX \right) \right] \\
%&=&\frac{1}{\sqrt{2m\hbar \omega }}\frac{1}{\sqrt{2m\hbar \omega }}\left[
%\left( P-im\omega X\right) ,\left( P+im\omega X\right) \right] \\
%&=&\frac{1}{2m\hbar \omega }\left\{ \left[ P,P\right] +\left[ P,im\omega X%
%\right] +\left[ -im\omega X,P\right] +\left[ -im\omega X,im\omega X\right] 
%\frac{{}}{{}}\right\} \\
%&=&\frac{1}{2m\hbar \omega }\left\{ im\omega \left[ P,X\right] -im\omega %
%\left[ X,P\right] -\left( im\omega \right) ^{2}\left[ X,X\right] \frac{{}}{{}%
%}\right\} \\
&=&\frac{1}{2m \hbar \omega }\left\{ -\ii m \omega \left[ \OX,\OP \right] -\ii m\omega \left[ \OX,\OP\right] \right\} =\frac{1}{2m\hbar \omega }\left(-2\ii m\omega \left[ \OX,\OP \right] \right) \\
&=&\frac{-\ii}{\hbar }\left[ \OX,\OP\right] =\frac{-\ii}{\hbar }\left( \ii\hbar\right) =-\ii^{2}=1.
\end{eqnarray*}

%\begin{equation*}
%\therefore ~\left[ \Oa,\Oa^{\dagger }\right] =1.
%\end{equation*}

\noindent
Ahora, al emplear (\ref{1.4}) en (\ref{1.8}):

\begin{eqnarray*}
\comm\big{\Oaa}{\Oaa} &=&\left[ \frac{1}{\sqrt{2m\hbar \omega }}\left( \OP-\ii m\omega \OX \right) ,\frac{1}{\sqrt{2m\hbar \omega }}\left( \OP-\ii m\omega \OX\right) \right] \\
&=&\frac{1}{2m\hbar \omega }\left( -\ii m\omega \left[ \OP,\OX \right] -\ii m \omega \left[ \OX,\OP \right] \right) \\
&=&\frac{1}{2m\hbar \omega }\left( \ii m\omega \left[ \OX,\OP \right] -\ii m\omega \left[ \OX,\OP \right] \right) =0.
\end{eqnarray*}

%\begin{equation*}
%\therefore ~ \comm\big{\Oa}{\Oa} =0.
%\end{equation*}

\noindent
Del mismo modo es posible comprobar que $\left[ \Oaa^{\dagger },\Oaa^{\dagger}\right] =0$.
\end{solution}


\noindent
Otra propiedad que se deriva fácilmente de las ecuaciones (\ref{1.3}) y (\ref{1.4}) es:

\begin{eqnarray}
\comm\bigg{\hat{H}}{\Oaa} &=&-\hbar \omega \Oaa  \label{1.9} \\
\left[ \hat{H},\Oaa^{\dagger }\right] &=&\hbar \omega \Oaa^{\dagger}  \label{1.9a}
\end{eqnarray}

\begin{exercise}
Demostrar los conmutadores (\ref{1.9}) y (\ref{1.9a})
\end{exercise}

\newcommand{\OoH}{\hat{H}}

\begin{solution}
Empleando la expresi\'{o}n para la Hamiltoniana (\ref{1.5}) y los resultados proporcionados en (\ref{1.7}) y (\ref{1.8}), el conmutador (\ref{1.9}) se puede calcular as\'{\i}:

\begin{eqnarray*}
\comm\bigg{\OoH}{\Oaa} &=&\left[ \frac{\hbar \omega }{2}\left( \Oaa^{\dagger} \Oaa+\Oaa \Oaa^{\dagger }\right) ,\Oaa \right] =\frac{\hbar \omega }{2}\left( \left[ \Oaa^{\dagger }\Oaa,\Oaa\right] +\left[ \Oaa \Oaa^{\dagger },\Oaa\right] \right) \\
%&=&\frac{\hbar \omega }{2}\left( a^{\dagger }\left[ a,a\right] +\left[
%a^{\dagger },a\right] a+a\left[ a^{\dagger },a\right] +\left[ a,a\right]
%a^{\dagger }\right) \\
&=&\frac{\hbar \omega }{2}\left( \left[ \Oaa^{\dagger },\Oaa \right] \Oaa + \Oaa \left[ \Oaa^{\dagger }, \Oaa \right] \right) =\frac{\hbar \omega }{2}\left( -\left[ \Oaa, \Oaa^{\dagger }\right] \Oaa - \Oaa \left[ \Oaa , \Oaa^{\dagger }\right] \right) \\
&=&\frac{\hbar \omega }{2}\left( -\Oaa-\Oaa\right) =\frac{-2\hbar \omega \Oaa}{2} =-\hbar \omega \Oaa,
\end{eqnarray*}

mientras que, para el conmutador (\ref{1.9a})

\begin{eqnarray*}
\comm\bigg{\OoH}{\Oaa^\dagger} &=&\left[ \frac{\hbar \omega }{2}\left( \Oaa^{\dagger }\Oaa+\Oaa \Oaa^{\dagger }\right) ,\Oaa^{\dagger }\right] =\frac{\hbar \omega }{2}\left( \left[ \Oaa^{\dagger }\Oaa,\Oaa^{\dagger }\right] +\left[ \Oaa \Oaa^{\dagger},\Oaa^{\dagger }\right]\right) \\
%&=&\frac{\hbar \omega }{2}\left( a^{\dagger }\left[ a,a^{\dagger }\right] +%
%\left[ a^{\dagger },a^{\dagger }\right] a+a\left[ a^{\dagger },a^{\dagger }%
%\right] +\left[ a,a^{\dagger }\right] a^{\dagger }\right) \\
&=&\frac{\hbar \omega }{2}\left( \Oaa^{\dagger }\left[ \Oaa,\Oaa^{\dagger }\right] +\left[ \Oaa,\Oaa^{\dagger }\right] \Oaa^{\dagger }\right) =\frac{\hbar \omega }{2}\left( \Oaa^{\dagger }\left[ \Oaa,\Oaa^{\dagger }\right] +\left[ \Oaa,\Oaa^{\dagger }\right] \Oaa^{\dagger }\right) \\
&=&\frac{\hbar \omega }{2}\left( \Oaa^{\dagger }+\Oaa^{\dagger }\right) =\frac{2\hbar \omega \Oaa^{\dagger }}{2}=\hbar \omega \Oaa^{\dagger }~.
\end{eqnarray*}

%concluyendo%
%\begin{equation*}
%\left[ H,a\right] =-\hbar \omega a\qquad ;\qquad \left[ H,a^{\dagger }\right]
%=\hbar \omega a^{\dagger }
%\end{equation*}
\end{solution}

\noindent
Por otro lado, observe que si se usa la relaci\'{o}n de conmutaci\'{o}n dada en (\ref{1.7}), se puede reescribir el hamiltoniano\ (\ref{1.5}) como:

\begin{eqnarray*}
\OoH &=&\frac{\hbar \omega }{2}\left[ \Oaa^{\dagger }\Oaa+\Oaa \Oaa^{\dagger }\right] =\frac{\hbar \omega }{2}\left[ \Oaa^{\dagger }\Oaa+\left( \Oaa^{\dagger } \Oaa + 1\right) \right]
\\
&=&\frac{\hbar \omega }{2}\left[ 2\Oaa^{\dagger }\Oaa+1\right]
\end{eqnarray*}%
donde se ha usado la relaci\'{o}n (\ref{1.7}), as\'{\i} se tiene,

\begin{equation}
\boxed{
\OoH=\hbar \omega \left( \Oaa^{\dagger} \Oaa+\frac{1}{2}\right)  .
}
\label{1.10}
\end{equation}

\noindent
Por otro lado, de los postulados de la mecanica cu\'{a}ntica se sabe que la ecuaci\'{o}n de Schr\"{o}dinger independiente del tiempo satisface la ecuaci\'{o}n de autovalores:

\begin{equation}
\OoH \left\vert n\right\rangle =E_{n}\left\vert n\right\rangle  \label{1.11}
\end{equation}

\noindent
donde $\left\vert n\right\rangle $ representa un auto-estado del operador Hamiltoniano con autovalor de energía $E_{n}$.
Es posible demostrar que $\Oaa \left\vert n\right\rangle $ y $\Oaa^{\dagger} \left\vert n\right\rangle $ forman un conjunto de autovectores asociados al operador $H$ con autovalores $\left( E_{n}-\hbar \omega \right) $ y $\left( E_{n}+\hbar \omega \right) $ respectivamente.

\vspace{0.3cm}
\begin{exercise}
Corroborar que $a\left\vert n\right\rangle $ y $a^{\dagger }\left\vert
n\right\rangle $ son autovectores del operador $H$ que tienen asociados los
autovalores $\left( E_{n}-\hbar \omega \right) $ y $\left( E_{n}+\hbar
\omega \right) $.
\end{exercise}
 
\begin{solution}

Para calcular la acci\'{o}n del operador $\OoH$ sobre $\Oaa |n\rangle $ es necesario analizar el conmutador (\ref{1.9}),

\begin{equation*}
\left[ \OoH,\Oaa \right] =\OoH \Oaa-\Oaa \OoH=-\hbar \omega \Oaa ~\ \implies ~\ \OoH \Oaa = \Oaa \OoH-\hbar \omega \Oaa,
\end{equation*}

de manera que 

\begin{eqnarray}
\OoH \left( \Oaa |n\rangle \right) &=&\left( \OoH \Oaa \right) \left\vert n\right\rangle =\left( \Oaa \OoH-\hbar \omega \Oaa \right) \left\vert n\right\rangle  \notag \\
&=& \Oaa \OoH\left\vert n\right\rangle -\hbar \omega \Oaa \left\vert n \right\rangle =\Oaa E_{n}\left\vert n\right\rangle -\hbar \omega \Oaa \left\vert n\right\rangle \notag
\end{eqnarray}

\begin{equation}
\boxed{
    \therefore ~ \OoH \left( \Oaa \left\vert n\right\rangle \right) =\left( E_{n}-\hbar \omega \right) \Oaa \left\vert n\right\rangle  \label{1.12a}.
    }
\end{equation}

en forma similar la acci\'{o}n de $\OoH$ sobre $\Oaa^{\dagger }\left\vert n\right\rangle $ se calcula al derivar de (\ref{1.9a}) el resultado,

\begin{equation*}
	\left[ \OoH,\Oaa^{\dagger }\right] =\OoH \Oaa^{\dagger }- \Oaa^{\dagger }\OoH=\hbar \omega \Oaa^{\dagger }~\ \implies ~\ \OoH \Oaa^{\dagger }=\Oaa^{\dagger } \OoH+\hbar \omega \Oaa^{\dagger},
\end{equation*}

así

\begin{eqnarray}
\OoH \left( \Oaa^{\dagger }|n\rangle \right) &=&\left( \OoH \Oaa^{\dagger }\right)
\left\vert n\right\rangle =\left( \Oaa^{\dagger } \OoH+\hbar \omega \Oaa^{\dagger
}\right) \left\vert n\right\rangle  \notag \\
&=& \Oaa^{\dagger } \OoH \left\vert n\right\rangle +\hbar \omega \Oaa^{\dagger
}\left\vert n\right\rangle =\Oaa^{\dagger }E_{n}\left\vert n\right\rangle
-\hbar \omega \Oaa^{\dagger }\left\vert n\right\rangle  \notag
\end{eqnarray}

\begin{equation}
    \boxed{
\therefore ~~ \OoH \left( \Oaa^{\dagger }|n \rangle \right) =\left( E_{n}+\hbar
\omega \right) \Oaa^{\dagger }|n\rangle  \label{1.12b}.
    }
\end{equation}

\end{solution}


De las expresiones dadas en (\ref{1.12a}) y (\ref{1.12b}) se observa que el operador \textquotedblleft $\Oaa$\textquotedblright\ parece aniquilar un cuanto de energ\'{\i}a, en la cantidad $\hbar \omega $ del estado considerado de energ\'{\i}a. Por otro lado, \textquotedblleft $\Oaa^{\dagger}$\textquotedblright\ crea o adiciona un cuanto de energía $\hbar \omega$ a $E_{n}$. Este sentido f\'{\i}sico permite identificarlos como operadores de aniquilaci\'{o}n ($\Oaa$) y creaci\'{o}n ($\Oaa^{\dagger}$), respectivamente. Por supuesto, en cualquier proceso f\'{\i}sico se debe conservar la energ\'{\i}a, por lo que los operadores $\Oaa$ o $\Oaa^{\dagger}$ no pueden aparecer solos en la expresi\'{o}n de ninguna cantidad mensurable. Por ejemplo, en la expresi\'{o}n del Hamiltoniano aparecen juntos, como se ve en la ecuaci\'{o}n. (\ref{1.9}).\\

\noindent
El estado fundamental o base del oscilador arm\'{o}nico se puede denotar mediante $\left\vert 0\right\rangle $. Dado que \'{e}ste es el estado de menor energ\'{\i}a, el operador de aniquilaci\'{o}n $\Oaa$, actuando sobre \'{e}l, no puede producir un estado de menor energ\'{\i}a. Por tanto, este estado debe ser totalmente aniquilado por la acci\'{o}n de $\Oaa$, lo que se puede escribir como:

\begin{equation}
	\Oaa \left\vert 0\right\rangle =0~\ \implies ~\text{Estado m\'{a}s bajo.}
	\label{1.13}
\end{equation}

Empleando la relaci\'{o}n (\ref{1.5}) es posible encontrar f\'{a}cilmente el autovalor de energ\'{\i}a asociado al estado fundamental, para ello hacemos

\begin{equation*}
\OoH \left\vert n\right\rangle =E_{n}\left\vert n\right\rangle ~\ \implies ~ 
\left[ \hbar \omega \left( \Oaa^{\dagger } \Oaa +\frac{1}{2}\right) \right]
\left\vert n\right\rangle =\hbar \omega \left( \Oaa^{\dagger } \Oaa \left\vert
n\right\rangle +\frac{1}{2}\left\vert n\right\rangle \right)
\end{equation*}
\begin{equation*}
\therefore ~\OoH\left\vert n\right\rangle =\hbar \omega \left( \Oaa^{\dagger
} \Oaa \left\vert n\right\rangle +\frac{1}{2}\left\vert n\right\rangle \right).
\end{equation*}

Al considerar diferentes valores para $n$,

\begin{itemize}
    \item $n=0$ 
\end{itemize}

\begin{equation*}
\OoH \left\vert 0\right\rangle =E_{0}\left\vert 0\right\rangle ~\ \implies ~ \OoH\left\vert 0\right\rangle =\hbar \omega \left( \Oaa^{\dagger }\underset{0}{\underbrace{ \Oaa \left\vert 0\right\rangle }}+\frac{1}{2}\left\vert
0\right\rangle \right) =E_{0}\left\vert 0\right\rangle
\end{equation*}


\begin{equation} \label{1.14}
\begin{split}
\frac{\hbar \omega }{2}\left\vert 0\right\rangle & = E_{0}\left\vert 0\right\rangle \\    
\therefore ~~ E_{0} &= \frac{\hbar \omega }{2},
\end{split}    
\end{equation}

Se debe tener en cuenta que en cu\'{a}ntica los valores que toma $n$ deben de ser netamente enteros. As\'{\i}, el primer estado excitado, que
denotaremos por $|1\rangle $, puede definirse por el estado cuya energ\'{\i}a es mayor que el estado fundamental en un cuanto, es decir,

\begin{equation}
|1\rangle =\Oaa^{\dagger }|0\rangle  \label{1.15},
\end{equation}

\newcommand{\OoN}{\hat{N}}

\begin{definition}[Operador Número]

Sea $\OoN$ un operador definido como:

\begin{equation*}
\OoN=\Oaa^{\dagger} \Oaa
\end{equation*}

a este se lo conoce como \emph{operador n\'{u}mero }y se cumple que la acci\'{o}n del mismo sobre el autoestado $\left\vert n\right\rangle $ satisface la relaci\'{o}n de autovalores

\begin{equation}
\OoN \left\vert n\right\rangle =n\left\vert n\right\rangle, 
\label{1.16}
\end{equation}
\label{def1}
\end{definition}


Con la definici\'{o}n (\ref{def1}) y empleando los operadores aniquilador y creador, se debe garantizar para un estado cu\'{a}ntico dado que

\begin{equation}
\Oaa \left\vert n\right\rangle =c_{n}\left\vert n-1\right\rangle ~\ \ \ \ ;~\ \
\ \ \Oaa^{\dagger }\left\vert n\right\rangle =c_{n}^{\prime }\left\vert
n+1\right\rangle   \label{1.15a}
\end{equation}

donde $c_{n}$ y $c_{n}^{\prime }$ se conocen como constantes de normalizaci\'{o}n. A continuaci\'{o}n se debe determinar el valor de dichas constantes, para ello se calcula el adjunto herm\'{\i}tico de las relaciones dadas en
(\ref{1.15a}), entonces:

\begin{eqnarray}
\Oaa \left\vert n\right\rangle  &=&c_{n}\left\vert n-1 \right\rangle \text{ \ \ \ 
}\overset{\dagger }{\longrightarrow }\text{ \ \ \ }\left\langle n\right\vert
\Oaa^{\dagger }=c_{n}^{\ast }\left\langle n-1\right\vert  \label{1.17}
\\
\Oaa^{\dagger }\left\vert n\right\rangle  &=&c_{n}^{\prime }\left\vert
n+1\right\rangle \text{ \ \ \ }\overset{\dagger }{\longrightarrow }\text{ \
\ \ }\left\langle n\right\vert \Oaa = c_{n}^{\prime \ast }\left\langle n+1\right\vert   \label{1.17a}
\end{eqnarray}

ahora, se calcula 

\begin{eqnarray}
\left\langle n\right\vert \Oaa \Oaa^{\dagger }\left\vert n\right\rangle &=&\left\langle n\right\vert 1+\Oaa^{\dagger } \Oaa \left\vert n\right\rangle =\langle n\left\vert n\right\rangle +\left\langle n\right\vert \left(a^{\dagger }a\left\vert n\right\rangle \right) =1+\left\langle n\right\vert \OoN \left\vert n\right\rangle =1+\left\langle n\right\vert n\left\vert n\right\rangle   \notag \\
&=&1+\langle n\left\vert n\right\rangle =1+n  \label{1.17b},
\end{eqnarray}

donde se ha empleado el conmutador (\ref{1.7}). Por otro lado es v\'{a}lido afirmar que  

\begin{eqnarray}
\langle n| \Oaa \Oaa^{\dagger }|n\rangle  &=&c_{n}^{\prime \ast }\langle
n+1|c_{n}^{\prime }|n+1\rangle =c_{n}^{\prime \ast }c_{n}^{\prime }\text{\ }%
\langle n+1|n+1\rangle   \notag \\
&=&|c_{n}^{\prime }|^{2}  \label{1.17c}.
\end{eqnarray}

Las expresiones (\ref{1.17b}) y (\ref{1.17c}) mostradas anteriormente, implican:

\begin{equation*}
|c_{n}^{\prime }|^{2}=n+1~\ \ \implies ~\ c_{n}^{\prime }=\sqrt{n+1}.
\end{equation*}

Entonces la acci\'{o}n del operador  $\Oaa^{\dagger }~$sobre$~|n\rangle $ es 

\begin{equation}
\boxed{
\Oaa^{\dagger }\left\vert n\right\rangle =\sqrt{n+1}\left\vert n+1\right\rangle .
\label{1.18}
}
\end{equation}

De la misma manera, se calcula

\begin{eqnarray}
\langle n|\Oaa^{\dagger } \Oaa \left\vert n\right\rangle  &=&\left\langle
n\right\vert \OoN \left\vert n\right\rangle =\left\langle n\right\vert
n\left\vert n\right\rangle   \notag \\
&=&n\langle n\left\vert n\right\rangle =n,  \label{1.18a}
\end{eqnarray}

y 

\begin{eqnarray}
\langle n|\Oaa^{\dagger } \Oaa \left\vert n\right\rangle  &=&c_{n}^{\ast }\langle
n-1|c_{n}\left\vert n-1\right\rangle \text{ }=c_{n}^{\ast }c_{n}\text{\ }%
\langle n-1\left\vert n-1\right\rangle   \notag \\
&=&|c_{n}|^{2}.  \label{1.18b}
\end{eqnarray}

resultados que permiten demostrar el resultado:

\begin{equation*}
|c_{n}|^{2}=n~\ \ \implies ~\ c_{n}=\sqrt{n}.
\end{equation*}

Por lo tanto, la acci\'{o}n del operador $\Oaa$ sobre $|n\rangle $ es
 
\begin{equation}
\boxed{
\Oaa \left\vert n\right\rangle =\sqrt{n}\left\vert n-1\right\rangle 
\label{1.19}.
}
\end{equation}

Se observa que del resultado (\ref{1.18}) se debe cumplir

\begin{equation}
\left\vert n+1\right\rangle =\frac{1}{\sqrt{n+1}} \Oaa^{\dagger }\left\vert
n\right\rangle   \label{1.20},
\end{equation}

pero al considerar diferentes valores enteros de $n$:

\begin{eqnarray*}
n &=&0\rightarrow |0+1\rangle =|1\rangle =\frac{1}{\sqrt{1}}\Oaa^{\dagger
}|0\rangle \, \\
n &=&1\rightarrow |1+1\rangle =|2\rangle =\frac{1}{\sqrt{2}}\Oaa^{\dagger
}|1\rangle =\frac{1}{\sqrt{2\cdot 1}}\left( \Oaa^{\dagger }\right)
^{2}|0\rangle \, \\
n &=&2\rightarrow |2+1\rangle =|3\rangle =\frac{1}{\sqrt{3}}\Oaa^{\dagger
}|2\rangle \,=\frac{1}{\sqrt{3\cdot 2\cdot 1}}\left( \Oaa^{\dagger }\right)
^{3}|0\rangle \, \\
%n &=&3\rightarrow |3+1\rangle =|4\rangle =\frac{1}{\sqrt{4}}a^{\dagger}|3\rangle \,=\frac{1}{\sqrt{4\cdot 3\cdot 2\cdot 1}}\left( a^{\dagger}\right) ^{4}|0\rangle \, \\
%n &=&4\rightarrow |4+1\rangle =|5\rangle =\frac{1}{\sqrt{5}}a^{\dagger}|4\rangle \,=\frac{1}{\sqrt{5\cdot 4\cdot 3\cdot 2\cdot 1}}\left(a^{\dagger }\right) ^{5}|0\rangle \,
\end{eqnarray*}

Generalizando, se puede concluir que:

\begin{equation}
\boxed{
|n\rangle =\frac{1}{\sqrt{n!}}\left( \Oaa^{\dagger }\right) ^{n}|0\rangle \
\label{1.21}.
}
\end{equation}

A la relaci\'{o}n anterior se la identifica como el \emph{n-\'{e}simo estado excitado}, mientras que $\frac{1}{\sqrt{n!}}$ representa la constante de
normalizaci\'{o}n, por lo tanto es un autoestado normalizado.

De lo anterior, tambi\'{e}n se puede determinar el autovalor de energ\'{\i}a para el oscilador armónico, as\'{\i}:

\begin{eqnarray*}
\OoH \left\vert n\right\rangle  &=&E_{n}\left\vert n\right\rangle ~\ \ \implies
~\ ~\left[ \hbar \omega \left( \Oaa^{\dagger } \Oaa+\frac{1}{2}\right) \right]
\left\vert n\right\rangle =\hbar \omega \left( \OoN|n\rangle +\frac{1}{2}%
\left\vert n\right\rangle \right) =E_{n}\left\vert n\right\rangle 
\end{eqnarray*}
\begin{equation}
    \therefore ~~ \OoH \left\vert n\right\rangle =\hbar \omega \left( n+\frac{1}{2}%
\right) \left\vert n\right\rangle =E_{n}\left\vert n\right\rangle, 
\end{equation}
 tal que

\begin{equation}
E_{n}=\hbar \omega \left( n+\frac{1}{2}\right)   \label{1.22}.
\end{equation}

\section{Relatividad Especial}
Dado que el objetivo es discutir las teor\'{\i}as de campo relativistas, aqu\'{\i} se presenta un resumen r\'{a}pido de la estructura matem\'{a}tica de la
relatividad especial, que ayudar\'{a} a establecer la notaci\'{o}n. En relatividad especial, las tres coordenadas espaciales y el tiempo definen
las cuatro componentes del cuadrivector de posici\'{o}n, que denotaremos por $x^{\mu }$. De acuerdo con la pr\'{a}ctica establecida, se utilizar\'{a}n 
\'{\i}ndices griegos $\mu ,\nu ,\lambda ,...$ para denotar los componentes de un cuadrivector, y los \'{\i}ndices latinos $i,j,k,...~$para denotar sus
componentes espaciales. En otras palabras, los \'{\i}ndices griegos toman valores $0,1,2,3$, mientras que los \'{\i}ndices latinos toman valores $1,2,3$. Por lo tanto,

\begin{equation}
x^{\mu }=(x^{0},x^{i})=(ct,\vb{x})  \label{18},
\end{equation}

donde se pone el factor de $c$ para que todos los componentes tengan dimensiones de longitud. Aqu\'{\i}, $c$ es la velocidad de la luz en el vac%
\'{\i}o, que es independiente del marco de acuerdo con uno de los axiomas fundamentales de la relatividad especial. La distancia $l$ entre dos puntos $%
x$ e $y$ en el espacio-tiempo se puede escribir en t\'{e}rminos de las coordenadas cartesianas como

\begin{equation}
l^{2}=(x^{0}-y^{0})^{2}-\sum\limits_{i=1}^{3}(x^{i}-y^{i})^{2}  \label{19}.
\end{equation}

Sin embargo, esta expresión se reescribe en forma compacta definiendo el tensor m\'{e}trico $g_{\mu \nu }$ de la siguiente manera:

\begin{equation}
l^{2}=g_{\mu \nu }(x^{\mu }-y^{\mu })(x^{\nu }-y^{\nu })  \label{20},
\end{equation}

donde las componentes del tensor m\'{e}trico son: 
\begin{equation}
g_{\mu \nu }=\left( 
\begin{array}{cccc}
1 & 0 & 0 & 0 \\ 
0 & -1 & 0 & 0 \\ 
0 & 0 & -1 & 0 \\ 
0 & 0 & 0 & -1%
\end{array}
\right)   \label{21}.
\end{equation}

Observe que en (\ref{20}) se ha utilizado la convenci\'{o}n de suma de Einstein, estableciendo que cualquier \'{\i}ndice que aparezca dos veces en el mismo t\'{e}rmino se
suma. No obstante, es posible usar la m\'{e}trica para bajar los \'{\i}ndices. Por ejemplo, es posible escribir un cuadrivector \emph{covariante }%
en la forma 

\begin{equation}
x_{\mu }=g_{\mu \nu }x^{\nu }=(ct,-\mathbf{x)}  \label{22}.
\end{equation}

De manera similar, se puede usar la inversa de la matriz $g_{\mu \nu }$, que
ser\'{a} denotada por $g^{\mu \nu }$, para elevar los \'{\i}ndices, es decir
se puede escribir un vector contravariante

\begin{equation}
x^{\mu }=g^{\mu \nu }x_{\nu }  \label{23}.
\end{equation}

Una propiedad importante es que la matriz inversa debe satisfacer la relaci\'{o}n

\begin{equation}
g^{\mu \nu }g_{\nu \lambda }=\delta _{\lambda }^{\mu }=g_{\lambda }^{\mu },
\label{24}
\end{equation}

donde $\delta _{\lambda }^{\mu }$ es el delta de kronecker en 4D, definido de tal modo que
\begin{equation}
\delta _{\lambda }^{\mu }=\left\{ 
\begin{array}{c}
1\text{ si }\mu =\lambda  \\ 
0\text{ si }\mu \neq \lambda 
\end{array}%
\right. \text{.}  \label{25}
\end{equation}

\begin{exercise}
Calcular la forma explícita de la inversa de la m\'{e}trica $g^{\mu \nu }$.
\end{exercise}
\begin{solution}
Para determinar la forma explicita de la métrica en nuestro caso, se hará uso del
siguiente teorema:
\begin{theorem}[Inversa de una matriz diagonal]
La matriz diagonal
\begin{equation*}
 D=\text{diag}(a_{1},a_{2},\dots,a_{n}),   
\end{equation*}

es invertible si y sólo si las entradas $a_{1},a_{2},\dots,a_{n}$ son todas distintas de $0$. En este caso, se tiene:
\begin{equation}
   D^{-1}=\text{diag}(a_{1}^{-1},a_{2}^{-1},\dots,a_{n}^{-1}) 
\end{equation}
\label{teorMaDiag}
\end{theorem}
Como la representación de la métrica dada en (\ref{21}) es una matriz diagonal, el teorema \ref{teorMaDiag} implica que la inversa de ésta sea de la forma:

\begin{equation}
g^{\mu \nu }=\left( 
\begin{array}{cccc}
\frac{1}{1} & 0 & 0 & 0 \\ 
0 & \frac{1}{-1} & 0 & 0 \\ 
0 & 0 & \frac{1}{-1} & 0 \\ 
0 & 0 & 0 & \frac{1}{-1}
\end{array}
\right)   \label{25a},
\end{equation}

\begin{comment}
Empleando la relaci\'{o}n (\ref{24}) y (\ref{21}): 
\begin{eqnarray*}
\left( 
\begin{array}{cccc}
a_{11} & a_{12} & a_{13} & a_{14} \\ 
a_{21} & a_{22} & a_{23} & a_{24} \\ 
a_{31} & a_{32} & a_{33} & a_{34} \\ 
a_{41} & a_{42} & a_{43} & a_{44}%
\end{array}%
\right) \left( 
\begin{array}{cccc}
1 & 0 & 0 & 0 \\ 
0 & -1 & 0 & 0 \\ 
0 & 0 & -1 & 0 \\ 
0 & 0 & 0 & -1%
\end{array}%
\right)  &=&\left( 
\begin{array}{cccc}
1 & 0 & 0 & 0 \\ 
0 & 1 & 0 & 0 \\ 
0 & 0 & 1 & 0 \\ 
0 & 0 & 0 & 1%
\end{array}%
\right)  \\
\left( 
\begin{array}{cccc}
a_{11} & -a_{12} & -a_{13} & -a_{14} \\ 
a_{21} & -a_{22} & -a_{23} & -a_{24} \\ 
a_{31} & -a_{32} & -a_{33} & -a_{34} \\ 
a_{41} & -a_{42} & -a_{43} & -a_{44}%
\end{array}%
\right)  &=&\left( 
\begin{array}{cccc}
1 & 0 & 0 & 0 \\ 
0 & 1 & 0 & 0 \\ 
0 & 0 & 1 & 0 \\ 
0 & 0 & 0 & 1%
\end{array}%
\right) 
\end{eqnarray*}%
la igualdad entre matrices se cumple si 
\begin{equation*}
a_{12}=a_{13}=a_{14}=a_{21}=a_{23}=a_{24}=a_{31}=a_{32}=a_{34}=a_{41}=a_{42}=a_{42}=0
\end{equation*}%
y%
\begin{equation*}
a_{11}=1~\ \ ;~\ \ a_{22}=a_{33}=a_{44}=-1
\end{equation*}
\end{comment}


lo que conlleva la forma explícita:
\begin{equation}
g^{\mu \nu }=\left( 
\begin{array}{cccc}
1 & 0 & 0 & 0 \\ 
0 & -1 & 0 & 0 \\ 
0 & 0 & -1 & 0 \\ 
0 & 0 & 0 & -1%
\end{array}
\right) =g_{\mu \nu }  \label{25b}.
\end{equation}
Este resultado indica que \textcolor{blue}{la inversa de la métrica coincide con la misma métrica}.
\end{solution}

Ahora, se estudia la transformación de un vector contravariante ante una transformación de Lorentz. Para ello se parte de la siguiente transformación:

\begin{equation}
x^{\mu}\rightarrow x^{\prime \mu}=\Lambda^{\mu}_{~\nu} x^{\nu},
\label{26}
\end{equation}

donde la forma espec\'{\i}fica de la matriz de transformación $\Lambda$ no es muy importante para nuestros prop\'{o}sitos. Cabe resaltar que estas transformaciones se definen
de modo que dejen invariante el producto escalar $x^{\mu }x_{\mu}$, entonces se debe garantizar que

\begin{eqnarray}
x^{\prime \mu }x_{\mu }^{\prime } &=&g_{\mu \nu }x^{\prime \mu }x^{\prime
\nu }=g_{\mu \nu }\Lambda_{~\rho}^{\mu}x^{\rho}\Lambda _{~\sigma }^{\nu
}x^{\sigma}  \label{27} \\
&=&g_{\mu \nu }\Lambda _{~\rho}^{\mu}\Lambda _{~\sigma}^{\nu}x^{\rho}x^{\sigma}  \notag,
\end{eqnarray}

sea a su ves igual a $x_{\sigma }x^{\sigma}$, que es equivalente a $g_{\rho \sigma}x^{\rho}x^{\sigma}$. Este planteamiento se cumple si

\begin{equation}
g_{\mu \nu }\Lambda _{~\rho}^{\mu }\Lambda _{~\sigma }^{\nu }=g_{\rho \sigma }.
\label{28}
\end{equation}

As\'{\i}, lo que las matrices de transformaci\'{o}n de Lorentz deben
satisfacer es que la m\'{e}trica no cambie cuando se pase de un sistema de
referencia a otro, el cual esta representado por la transformaci\'{o}n de
coordenadas mediadas por las transformaciones de Lorentz.

\begin{exercise}
Analizar la forma en cómo transforma un cuadrivector covariante.
\end{exercise}
\begin{solution}
Una cantidad covariante mediada bajo una transformaci\'{o}n de Lorentz se
deber\'{\i}a escribir como
\begin{equation}
x_{\mu}^{\prime}=\Lambda _{\mu}^{~\nu}x_{\nu}   \label{29}.
\end{equation}
Con el f\'{\i}n de demostrar esto, si se utiliza la relaci\'{o}n (\ref{26})
y se le multiplica al lado izquierdo de ambos lados por un componente de la m\'{e}trica, 
es decir
\begin{eqnarray*}
x^{\prime \mu } &=&\Lambda_{~\nu}^{\mu}x^{~\nu}~\ \ \implies ~\ \ g_{\rho
\mu }x^{\prime \mu }=g_{\rho \mu }\Lambda _{~\nu}^{\mu}x^{\nu}=\Lambda
_{\rho \nu }x^{\nu } \\
&\therefore &~x_{\rho }^{\prime }=\Lambda_{\rho \nu }x^{\nu },
\end{eqnarray*}
y usando nuevamente la m\'{e}trica en el lado derecho de la línea anterior
para bajar el \'{\i}ndice $\nu $
\begin{equation*}
~x_{\rho }^{\prime}=\Lambda_{\rho \nu}g^{\nu \sigma}x_{\sigma}=\Lambda_{\rho}^{~\sigma}x_{\sigma }.
\end{equation*}

por lo tanto se deduce que
\begin{equation}
x_{\rho }^{\prime }=\Lambda _{\rho }^{~\sigma }x_{\sigma }  \label{29b}.
\end{equation}

\end{solution}

\begin{exercise}
Demostrar que $A^{\mu }B_{\mu}=A_{\mu}B^{\mu}$
\end{exercise}
\begin{solution}
La manera más sencilla de mostrar lo planteado es haciendo uso de la métrica para subir y bajar índices, es decir:
\begin{equation}
    A^{\mu}B_{\mu}=A^{\mu}g_{\mu \nu}B^{\nu}=A_{\nu}B^{\nu}
    \label{30}.
\end{equation}
\begin{comment}
Escribiendo los cuadrivectores en la forma (\ref{18}) 
\begin{eqnarray*}
A^{\mu }B_{\mu } &=&A_{\mu }B^{\mu } \\
\left( A^{0},A^{i}\right) \cdot \left( B_{0},B_{i}\right)  &=&\left(
A_{0},A_{i}\right) \cdot \left( B^{0},B^{i}\right) 
\end{eqnarray*}%
usando la m\'{e}trica para subir el \'{\i}ndice espacial y temporal%
\begin{equation*}
\left( A^{0},A^{i}\right) \cdot \left( g^{00}B_{0},g_{ij}B^{j}\right)
=\left( g_{00}A^{0},g_{ij}A^{j}\right) \cdot \left( B^{0},B^{i}\right) 
\end{equation*}%
pero para las componentes $g_{ij}$ de la m\'{e}trica, los \'{u}nicos
valores diferentes de cero de acuerdo a la forma explicita (\ref{21}), son
los diagonales es decir cuando $j=i~\ $ y se cumple $g_{ii}=$ $-1$ en el
caso espacial, mientras que en el caso temporal $g_{00}=1$. En esa forma  
\begin{eqnarray}
\left( A^{0},A^{i}\right) \cdot \left( g^{00}B_{0},g_{ii}B^{i}\right) 
&=&\left( g_{00}A^{0},g_{ii}A^{i}\right) \cdot \left( B^{0},B^{i}\right)  
\notag \\
\left( A^{0},A^{i}\right) \cdot \left( B^{0},-B^{i}\right)  &=&\left(
A^{0},-A^{i}\right) \cdot \left( B^{0},B^{i}\right)   \notag \\
A^{0}B^{0}-A^{i}B^{i} &=&A^{0}B^{0}-A^{i}B^{i}~~\impliedby \text{Invariante
de Minkowsky}  \notag \\
A^{\mu }B_{\mu } &=&A_{\mu }B^{\mu }  \label{29a}
\end{eqnarray}%
Otra forma de verificar lo anterior es usando la m\'{e}trica, es decir%
\begin{remark}
\begin{eqnarray}
A^{\mu }B_{\mu } &=&A^{\mu }g_{\mu \nu }B^{\nu }=A_{\nu }B^{\nu }  \notag \\
&=&\left( 
\begin{array}{cccc}
A^{0} & A^{1} & A^{2} & A^{3}%
\end{array}%
\right) \left( 
\begin{array}{cccc}
1 & 0 & 0 & 0 \\ 
0 & -1 & 0 & 0 \\ 
0 & 0 & -1 & 0 \\ 
0 & 0 & 0 & -1%
\end{array}%
\right) \left( 
\begin{array}{c}
B^{0} \\ 
B^{1} \\ 
B^{2} \\ 
B^{3}%
\end{array}%
\right)   \notag \\
&=&\left( 
\begin{array}{cccc}
A^{0} & -A^{1} & -A^{2} & -A^{3}%
\end{array}%
\right) \left( 
\begin{array}{c}
B^{0} \\ 
B^{1} \\ 
B^{2} \\ 
B^{3}%
\end{array}%
\right)   \notag \\
&=&A^{0}B^{0}-A^{1}B^{1}-A^{2}B^{2}-A^{3}B^{3}  \notag \\
&=&A^{0}B^{0}-A^{i}B^{i}  \label{30}
\end{eqnarray}
\end{remark}
\end{comment}
\end{solution}

\vspace{1cm}

\begin{cabox}
\emph{Cualquier objeto que permanece invariante como $x^{\mu }x_{\mu }$ representa una 
cantidad escalar. Además, un elemento de cuatro componentes que se transforma
como \textcolor{red}{$x^{\mu}$} se llama vector contravariante, mientras que un objeto de
cuatro componentes que se transforma como \textcolor{red}{$x_{\mu}$} es un vector covariante.}
\end{cabox}
\noindent

El producto escalar de dos cuadrivectores $A^{\mu }$ y $B^{\mu }$, con base en la relaci\'{o}n a (\ref{30}) se escribe como:

\begin{eqnarray}
A^{\mu }B_{\mu } &=&A_{\mu }B^{\mu }=g_{\mu \nu }A^{\mu }B^{\nu }  \label{31}
\\
&=&A^{0}B_{0}+A^{i}B_{i}  \notag \\
&=&A^{0}B^{0}-\mathbf{A}\cdot \mathbf{B}  \notag,
\end{eqnarray}

donde $\mathbf{A}\cdot \mathbf{B}$ es el producto escalar tridimensional
habitual. \\

Si $\phi (\mathbf{x})$ es una funci\'{o}n escalar multivarible, en
4D se reescribe esta funci\'{o}n escalar como $\phi (x^{\mu })$. Al considerar
un \emph{peque\~{n}o cambio infinitesimal en las coordenadas}


\begin{equation}
d\phi (x)=\frac{\partial \phi }{\partial x^{\mu }}dx^{\mu},
\label{32}
\end{equation}

se tiene que cumplir que  el lado derecho de (\ref{32}) sea un escalar. Pero $dx^{\mu }$\ es
un vector contravariante, por lo que $\frac{\partial \phi }{dx^{\mu}}$ deberá ser un vector covariante, que 
de ahora en adelante se denotar\'{a} por $\partial _{\mu }\phi 
$. Con esta identificaci\'{o}n:

\begin{eqnarray}
\partial _{\mu } &\equiv &\frac{\partial }{\partial x^{\mu }}=\left( \frac{1%
}{c}\frac{\partial }{\partial t},\frac{\partial }{\partial x},\frac{\partial 
}{\partial y},\frac{\partial }{\partial z}\right)   \label{33} \\
&=&\left( \frac{1}{c}\frac{\partial }{\partial t},\nabla \right).  \notag
\end{eqnarray}

Tenga presente que dado un vector arbitrario $A^{\mu }$ se garantizará que

\begin{equation}
\partial _{\mu }A^{\mu }=\frac{1}{c}\frac{\partial A^{0}}{\partial t}+\nabla
\cdot \mathbf{A}  \label{34}.
\end{equation}

Otro operador a considerar es el \emph{operador de onda}, que representa un 
operador escalar escrito como,
\begin{equation}
\square \equiv \frac{1}{c^{2}}\frac{\partial ^{2}}{\partial t^{2}}-\nabla
^{2}=\partial _{\mu }\partial ^{\mu}.  \label{35}
\end{equation}

Por otro lado, recuerde que el cuadrimomento de una part\'{\i}cula es dado por:

\begin{equation}
p^{\mu }=\left( \frac{E}{c},\vb{p}\right),  \label{36}
\end{equation}

donde $E$ es la energ\'{\i}a de la part\'{\i}cula y $\textbf{p}$ su momento lineal asociado. El producto escalar
invariante $p^{\mu }p_{\mu }$ es dado por:

\begin{eqnarray}
p^{\mu }p_{\mu } &=&p^{0}p^{0}-\mathbf{p}^{2}  \notag \\
&=&\frac{E^{2}}{c^{2}}-\mathbf{p}^{2}  \notag \\
&=&m^{2}c^{2}  \label{37a},
\end{eqnarray}%

donde (\ref{37a}) viene de la definici\'{o}n de la cuadrivelocidad $\eta
_{\mu }\eta ^{\mu }=c^{2}$ que es una cantidad invariante y $m$ es la masa
en reposo de la part\'{\i}cula, que llamaremos simplemente masa.

\section{Espacio y tiempo en teor\'{\i}a cu\'{a}ntica relativista}

??

\begin{comment}
En la mec\'{a}nica cu\'{a}ntica no relativista, el estado del tiempo y de
las coordenadas espaciales son muy diferentes, como se menciono
anteriormente. En una teor\'{\i}a relativista, deben equipararse antes de
continuar. Esto se logra tratando el espacio en la misma base que el tiempo,
es decir, tratando las coordenadas espaciales tambi\'{e}n como par\'{a}%
metros. As\'{\i}, las cantidades de las que se puede hablar en las teor\'{\i}%
as cu\'{a}nticas relativistas son funciones tanto del espacio como del
tiempo, es decir, del espacio-tiempo. Estas funciones se denominan \emph{%
campos} incluso en la\textbf{\ f\'{\i}sica cl\'{a}sica}. Por ejemplo, en la f%
\'{\i}sica cl\'{a}sica, cuando se habla del campo electromagn\'{e}tico, se
relaciona el campo el\'{e}ctrico y el campo magn\'{e}tico que son funciones
de las coordenadas espaciales y del tiempo. Del mismo modo, en la teor\'{\i}%
a cu\'{a}ntica de campos, se hablar\'{a} de varios tipos de campos, incluido
el campo electromagn\'{e}tico. En mec\'{a}nica cl\'{a}sica o cu\'{a}ntica,
se podr\'{\i}a hablar de objetos que eran funciones del tiempo \'{u}%
nicamente. Pero la necesidad de hablar de espacio y tiempo en pie de
igualdad requiere autom\'{a}ticamente que se hable de campos.

Es importante darse cuenta de que las relaciones de incertidumbre, que
forman la base de una teor\'{\i}a cu\'{a}ntica, deben interpretarse de
manera diferente para este prop\'{o}sito. En la mec\'{a}nica cu\'{a}ntica no
relativista, la relaci\'{o}n

\begin{equation}
\Delta x\Delta p_{x}\gtrsim \hbar   \label{38}
\end{equation}

se interpreta en t\'{e}rminos de las incertidumbres en las medidas de la
coordenada $x$ y el momento $p_{x}$ correspondiente. Por otro lado, la relaci%
\'{o}n de incertidumbre tiempo-energ\'{\i}a,

\begin{equation}
\Delta t\Delta E\gtrsim \hbar   \label{39}
\end{equation}

se interpreta diciendo que si se intenta realizar una medici\'{o}n de energ%
\'{\i}a con una precisi\'{o}n $\Delta E$, el proceso de medici\'{o}n deber%
\'{\i}a llevar un tiempo $\Delta t$ que est\'{e} relacionado con $\Delta E$
por la relación anterior. Ahora que $x$ tambi\'{e}n es un par\'{a}metro,
la interpretaci\'{o}n de la ecuaci\'{o}n (\ref{38}) tambi\'{e}n debe seguir
la l\'{\i}nea de la de la ecuaci\'{o}n (\ref{39}). En otras palabras,
debemos entender que la medici\'{o}n del momento con una precisi\'{o}n $%
\Delta p_{x}$ requiere que la medici\'{o}n se realice en una regi\'{o}n cuya
extensi\'{o}n espacial sea mayor que $\Delta x$ dada por la ecuaci\'{o}n (%
\ref{38}).
\end{comment}
\begin{comment}
En la mecánica cuántica no relativista, el estado del tiempo y de las coordenadas espaciales son
muy diferentes, como se menciono anteriormente. En una teoría relativista, deben equipararse antes
de continuar, esto se logra tratando el espacio en la misma base que el tiempo, es decir, considerando las
coordenadas espaciales también como parámetros. Así, las cantidades de las que se puede hablar
en las teorías cuánticas relativistas son funciones tanto del espacio como del tiempo (espacio-tiempo). Estas funciones se denominan \emph{campos} incluso en la física clásica.\\

A nivel cuántico se dieron dos hallazgos experimentales muy importantes como lo son los experimentos de Davisson-Germer y de doble rendija los cuales han demostrado que las partículas microscópico dan lugar a patrones de interferencia. Para explicar el patrón de interferencia, se vio la necesidad de describir las partículas microscópicas por medio de ondas. Las ondas no están localizadas en el espacio. Como resultado, tenemos que renunciar a la precisión para describir partículas microscópicas, ya que las ondas dan, en el mejor de los casos, una explicación probabilística. Por otro lado, se ha visto en el experimento de la doble rendija que es imposible seguir el movimiento de los electrones individuales; no existe ningún dispositivo experimental que pueda determinar la rendija a través de la cual ha pasado un electrón dado. No ser capaz de predecir eventos individuales es una flagrante violación de un principio fundamental de la física clásica: \textcolor{blue}{previsibilidad o determinación}. Estos hallazgos experimentales inspiraron a Heisenberg a postular la naturaleza indeterminista del mundo micro-físico y a Born a introducir la interpretación probabilística de la mecánica cuántica.

\subsection{Principio de incertidumbre de Heisenberg}

Según la física clásica, dadas las condiciones iniciales y las fuerzas que actúan sobre un sistema, se puede determinar con exactitud el comportamiento futuro (camino único) de este sistema físico. Es decir, si se conocen las coordenadas iniciales $\mathbf{r}_{0}$, la velocidad $\mathbf{v}_{0}$ y todas las fuerzas que actúan sobre la partícula, la posición $\mathbf{r}(t)$ y la velocidad $\mathbf{v}(t)$ se determinan de manera única mediante la segunda ley de Newton. La física clásica es pues, completamente \emph{determinista}.\\

¿Esta visión determinista es válida también para el mundo micro físico (cuántico)? Dado que una partícula se representa en el contexto de la mecánica cuántica por medio de una función de onda correspondiente a la onda de la partícula, y dado que las funciones de onda no se pueden localizar, entonces una partícula microscópica está algo dispersa espacio y, a diferencia de las partículas clásicas, no se puede localizar en el espacio. Además, hemos visto en el experimento de la doble rendija que es imposible determinar la rendija por la que pasó el electrón sin perturbarlo. Los conceptos clásicos de posición exacta, momento exacto y trayectoria única de una partícula, por lo tanto, no tienen sentido a escala microscópica. Esta es la esencia del principio de incertidumbre de Heisenberg.\\

En su forma original, el principio de incertidumbre de Heisenberg establece que:
\begin{cabox}
\emph{Si la componente $x$ del momento de una partícula se mide con una incertidumbre $\Delta p_{x}$, entonces su posición $x$ no se puede al mismo tiempo medirse con mayor precisión que, $\Delta x=\hbar/\left(2\Delta p_x\right)$.\\
La forma tridimensional de las relaciones de incertidumbre para la posición y el momento se puede escribir de la siguiente manera}:
\begin{equation}
    \Delta x\Delta p_x\geq\frac{\hbar}{2}\ \ \ \ ,\ \ \ \ \Delta y\Delta p_y\geq\frac{\hbar}{2}\ \ \ \ ,\ \ \ \ \Delta z\Delta p_z\geq\frac{\hbar}{2}.
\end{equation}
\end{cabox}
Este principio indica que, aunque es posible medir con precisión el momento o la posición de una partícula, no es posible medir estos dos observables simultáneamente con una precisión arbitraria. Es decir, no podemos localizar una partícula microscópica sin darle un momentum bastante grande. No podemos medir la posición sin perturbarla; no hay forma de llevar a cabo una medición de este tipo de forma pasiva, ya que está obligada a cambiar el impulso. Para comprender esto, considere medir la posición de un objeto macroscópico (p. ej., un automóvil) y la posición de un sistema microscópico (p. ej., un electrón en un átomo). Por un lado, para localizar la posición de un objeto macroscópico basta con observarlo; la luz que lo golpea y se refleja en el detector (sus ojos o un dispositivo de medición) no puede afectar de manera medible el movimiento del objeto. Por otro lado, para medir la posición de un electrón en un átomo, se debe usar la radiación de longitud de onda muy corta (del tamaño del átomo). La energía de esta radiación es lo suficientemente alta como para cambiar tremendamente el momento del electrón; la mera observación del electrón afecta tanto su movimiento que puede sacarlo completamente de su órbita. Por lo tanto, es imposible determinar la posición y el momento simultáneamente con una precisión arbitraria. Si se localizara una partícula, su función de onda sería cero en cualquier otro lugar y su onda tendría entonces una longitud de onda muy corta. \\

Según la relación de D’ Broglie $p=h/\lambda$, el momento de esta partícula será bastante alto. Formalmente, esto significa que si una partícula se localiza con precisión (es decir, $\Delta x\ \rightarrow\ 0$), habrá una incertidumbre total sobre su momento (es decir, $\Delta p_x\rightarrow\ \infty $). Para resumir, dado que todos los fenómenos cuánticos se describen mediante ondas, no tenemos más remedio que aceptar límites en nuestra capacidad para medir simultáneamente dos variables complementarias cualesquiera.\\

El principio de incertidumbre de Heisenberg se puede generalizar a cualquier par de variables dinámicas complementarias o canónicamente conjugadas: es imposible diseñar un experimento que pueda medir simultáneamente dos variables complementarias con una precisión arbitraria (si esto se lograra alguna vez, la teoría de la mecánica cuántica colapsaría).\\

La energía y el tiempo, por ejemplo, forman un par de variables complementarias. Su medida simultánea debe obedecer a la relación de incertidumbre tiempo-energía:
\begin{equation}
    \Delta E\Delta t\geq\frac{\hbar}{2}.
\end{equation}

Esta relación establece que, si hacemos dos medidas de la energía de un sistema y si estas medidas están separadas por un intervalo de tiempo $\Delta t$, las energías medidas diferirán en una cantidad $\Delta E$ que en ningún caso puede ser menor que $\hbar/\Delta t$. Si el intervalo de tiempo entre las dos mediciones es grande, la diferencia de energía será pequeña. Esto se puede atribuir a que, cuando se realiza la primera medición, el sistema se perturba y tarda mucho en volver a su estado inicial sin perturbaciones. Esta expresión es particularmente útil en el estudio de los procesos de descomposición, ya que especifica la relación entre el tiempo de vida media y el ancho de energía de los estados excitados.\\

Vemos que, en marcado contraste con la física clásica, la mecánica cuántica es una teoría completamente indeterminista. Al preguntar sobre la posición o el momento de un electrón, no se puede obtener una respuesta definitiva; sólo es posible una respuesta probabilística. Según el principio de incertidumbre, si la posición de un sistema cuántico está bien definida, su cantidad de movimiento será totalmente indefinida. En este contexto, el principio de incertidumbre ha derribado claramente uno de los conceptos más importantes de la física clásica: el carácter determinista de la mecánica newtoniana.
\end{comment}
\section{Unidades Naturales}

Los cursos elementales de f\'{\i}sica comienzan con tres unidades
independientes, longitud, tiempo y masa. A medida que se avanza, se
introducen unidades adicionales, por ejemplo, la carga el\'{e}ctrica,
temperatura, etc.

En las discusiones sobre la teor\'{\i}a cu\'{a}ntica de campos, se
acostumbra utilizar un sistema de unidades en el que s\'{o}lo hay una unidad
fundamental, que se puede considerar como la unidad de masa. Las unidades de
longitud y tiempo se definen declarando que
\begin{equation}
\hbar =1\qquad \text{y}\qquad c=1  \label{40}.
\end{equation}
Las unidades resultantes se denominan \emph{unidades naturales}, que se
utilizar\'{a}n en el resto de estos apuntes. Las dimensiones de varias
cantidades en estas unidades se presentan en la tabla (\ref{tab:table1}). Hemos
dado los factores de conversi\'{o}n para cantidades expresadas en unidades
de energ\'{\i}a, que es universalmente preferido en los c\'{a}lculos. Se
deben tener en cuenta que la mayor\'{\i}a de los otros textos dan factores
de conversi\'{o}n para unidades de masa.

\begin{table}[h!]
\renewcommand{\tablename}{Tabla}
\begin{center}
\begin{tabular}{ c | c | c }
\hline \hline
Cantidad & Dimensión &  factor de conversión\\ \hline \hline
Masa & $\left[M \right]$ & $1/c^{2}$\\
Longitud & $\left[M \right]^{-1}$ & $\hbar c$\\
Tiempo & $\left[M \right]^{-1}$ & $\hbar$\\ 
Energía & $\left[M \right]$ & $1$\\	
Momento & $\left[M \right]$ & $1/c$\\
Fuerza & $\left[M \right]^{2}$ & $1/\hbar c$\\
Acción & $\left[M \right]^{0}$ & $\hbar$\\
Carga eléctrica & $\left[M \right]^{0}$ & $\sqrt{\hbar c}$\\ \hline \hline
\end{tabular}
\caption{Dimensiones de diversas cantidades físicas en unidades naturales.}
\label{tab:table1}
\end{center}
\end{table}

La carga el\'{e}ctrica se define de tal manera que la ley de Coulomb acerca
de la fuerza $\mathbf{F}$ entre dos cargas $q_{1}$ y $q_{2}$ toma la forma:

\begin{equation}
\mathbf{F}=\frac{q_{1}q_{2}}{4\pi r^{2}}\hat{r}=\frac{q_{1}q_{2}}{4\pi r^{3}}%
\mathbf{r~,}  \label{41}
\end{equation}

Por tanto, la carga el\'{e}ctrica no tiene dimensiones, como se muestra en
la Tabla \ref{tab:table1}. Se tomar\'{a} como unidad de carga el\'{e}ctrica la
carga del prot\'{o}n, $e$. Finalmente la constante de estructura fina $%
\alpha $ est\'{a} relacionada con esta unidad por la siguiente relaci\'{o}n:

\begin{equation}
\alpha =\frac{e^{2}}{4\pi },\qquad e>0\text{.}  \label{42}
\end{equation}

\begin{exercise}
En unidades naturales, el tiempo de vida a la inversa del mu\'{o}n viene dado por
\begin{equation}
\tau ^{-1}=\frac{G_{F}^{2}m^{5}}{192\pi ^{3}}~\text{,}  \label{43}
\end{equation}%
donde $m$ es la masa del muon, que vale $106$ MeV. \textquestiondown Cu\'{a}%
l es la dimensi\'{o}n de $G_{F}$ en unidades naturales? Introduzca los
factores de $\hbar $ y $c$ para que la ecuaci\'{o}n tambi\'{e}n se pueda
interpretar en unidades convencionales. A partir de esto, calcule el tiempo
de vida en segundos si $G_{F}=1,166\times 10^{-11}$ en unidades de MeV.
\end{exercise}
\begin{solution}
En unidades naturales $\hbar =1$  y $c=1$, lo que da como resultado que cada
cantidad f\'{\i}sica se exprese en unidades de masa. Como ejemplo, la vida 
\'{u}til del mu\'{o}n $\tau $ viene dada por 
\begin{equation*}
\tau ^{-1}=\frac{G_{F}^{2}m^{5}}{192\pi ^{3}},
\end{equation*}%
donde $m$ ya se conoce y $G_{F}$ es la \emph{constante de acoplamiento de
Fermi}, que surge de la teor\'{\i}a cu\'{a}ntica de campo de la interacci%
\'{o}n electrod\'{e}bil. Dado que el tiempo tiene unidades naturales de masa
inversa $M^{-1}$, las unidades de $G_{F}$ deben ser $M^{-2}$ para que la
igualdad quede dimensionalmente correcta.

Para convertir esta f\'{o}rmula en unidades del sistema internacional (SI),
necesitamos insertar los factores de $\hbar $ y $c$ para que el lado
izquierdo de la ecuaci\'{o}n (\ref{43}) tenga unidades de $s^{-1}$. Dado que
las unidades son actualmente $M^{1}$ donde la masa se expresa en MeV, que es
una unidad de energ\'{\i}a, y las unidades de $\hbar $ son de acci\'{o}n,
que es \emph{(energ\'{\i}a)}$\emph{\ }\times \emph{\ }$\emph{(tiempo)},
podemos dividir por $\hbar $ para obtener las unidades de \emph{(tiempo)}$%
^{-1}$. Por lo tanto%
\begin{equation*}
\tau ^{-1}=\frac{1}{\hbar }\frac{G_{F}^{2}m^{5}}{192\pi ^{3}}
\end{equation*}%
Dado que $G_{F}=1,166\times 10^{-11}~$MeV$^{-2}$~y $\hbar =6,58\times
10^{-22}$MeV$\cdot $s, obtenemos%
\begin{eqnarray}
\tau ^{-1} &=&\frac{1}{\left( 6,58\times 10^{-22}~\text{MeV}\cdot \text{s}%
\right) }\frac{\left( 1,166\times 10^{-11}~\text{MeV}^{-2}\right) ^{2}\left(
106~\text{MeV}\right) ^{5}}{192\pi ^{3}}  \notag \\
%&=&\frac{\left( 1,166\right) ^{2}\cdot ~\left( 106\right) ^{5}\times10^{-22}~\text{MeV}^{-4}~\text{MeV}^{5}}{6,58\cdot \left( 192\pi ^{3}\right)\times 10^{-22}~\text{MeV}\cdot \text{s}}  \notag \\
%&=&\frac{\left( 1,166\right) ^{2}\cdot ~\left( 106\right) ^{5}~\text{MeV}~}{%6,58\cdot \left( 192\pi ^{3}\right) ~\text{MeV}\cdot \text{s}}=\frac{\left(1,166\right) ^{2}\cdot ~\left( 106\right) ^{5}~~}{6,58\cdot \left( 192\pi^{3}\right) ~}\text{s}^{-1}  \notag \\
&=&4,645\times 10^{5}\text{s}^{-1}  \label{43a}
\end{eqnarray}%

es decir
\begin{equation*}
\tau =2,15\times 10^{-6}\text{s}.
\end{equation*}
\end{solution}

%-----------------------------------------------      Fin Capítulo 1 -------------------------------------------------------------------------------------------